\begingroup
\color{borrador2}
\textcolor{borrador2}{La siguiente figura propone una vista de linaje que incluye los siete orígenes de fragmento comprobados en el corpus completo.}

\begin{figure}[H]
  \centering
  \resizebox{0.98\textwidth}{!}{%
    \begin{tikzpicture}[
      nodo/.style={draw=borrador2, thick, rounded corners=3pt, fill=white,
                   text=borrador2, align=center, font=\small,
                   minimum width=2.7cm, minimum height=0.9cm},
      origen/.style={nodo, fill=borrador2!7, minimum width=3.2cm,
                     font=\footnotesize\ttfamily},
      central/.style={draw=borrador2, very thick, rounded corners=4pt,
                      fill=borrador2!10, text=borrador2, align=left,
                      font=\footnotesize, inner sep=7pt, text width=5.0cm},
      grupo/.style={draw=borrador2, densely dashed, rounded corners=5pt,
                    inner sep=9pt},
      titulo/.style={font=\small\bfseries, text=borrador2},
      flecha/.style={-{Latex[length=2.2mm]}, thick, draw=borrador2},
    ]
      \node[nodo] (grado) at (0,2.4) {Grado};
      \node[nodo] (asig) at (0,0.9) {Asignatura};
      \node[nodo] (guia) at (0,-0.6) {Guía docente};
      \node[nodo] (salidas) at (0,-2.1) {Salidas profesionales};

      \node[origen] (o1) at (5.0,2.8) {guia};
      \node[origen] (o2) at (5.0,1.7) {asignatura\_sin\_guia};
      \node[origen] (o3) at (5.0,0.6) {salidas};
      \node[origen] (o4) at (5.0,-0.5) {plan\_de\_estudios};
      \node[origen] (o5) at (5.0,-1.6) {catalogo};
      \node[origen] (o6) at (5.0,-2.7) {ficha\_titulacion};
      \node[origen] (o7) at (5.0,-3.8) {mencion};

      \node[central] (chunk) at (11.2,-0.5)
        {\textbf{Fragmento}\\[2pt]
         \texttt{origen}\\
         \texttt{grados[]} $\leftrightarrow$ \texttt{codigos[]}\\
         \texttt{nombre}, \texttt{texto}\\
         \texttt{chunk\_index}, \texttt{total\_chunks}};

      \draw[flecha] (grado.east) -- (o4.west);
      \draw[flecha] (grado.east) -- (o5.west);
      \draw[flecha] (grado.east) -- (o6.west);
      \draw[flecha] (grado.east) -- (o7.west);
      \draw[flecha] (asig.east) -- (o2.west);
      \draw[flecha] (guia.east) -- (o1.west);
      \draw[flecha] (salidas.east) -- (o3.west);

      \foreach \n in {o1,o2,o3,o4,o5,o6,o7}
        \draw[flecha] (\n.east) -- (chunk.west);

      \begin{scope}[on background layer]
        \node[grupo, fit=(grado)(asig)(guia)(salidas),
              label={[titulo]north:Registros extraídos}] {};
        \node[grupo, fit=(o1)(o2)(o3)(o4)(o5)(o6)(o7),
              label={[titulo]north:Orígenes generados}] {};
      \end{scope}

      \node[font=\scriptsize, text=borrador2, align=center, anchor=north]
        at (11.2,-2.7)
        {Las listas de titulaciones y códigos son paralelas.\\Una guía compartida puede pertenecer a varios grados.};
    \end{tikzpicture}%
  }
  \caption[Alternativa del modelo de datos]{\textcolor{borrador2}{Versión alternativa del modelo de datos orientada al linaje. Presenta los cuatro tipos extraídos, los siete valores reales del campo \texttt{origen} y la estructura común de los 1499 fragmentos verificados en el corpus del 19/08/2026. Fuente: Elaboración propia.}}
  \label{fig:modelo_datos_alt}
\end{figure}
\endgroup
